\section{Le point de vue de Copenhague}
\subsection{La description scientifique (Bohr, Heisenberg)}
La description de la Mécanique Quantique (MQ), que nous dénommerons plutôt
« description standard » vu les diverses nuances existant parmi les auteurs sur le point de
vue de Copenhague, est la suivante. De manière générale, le ‘vecteur d’état’ est associé à la
préparation [conditionnement à l’équilibre thermique par exemple] d’un système physique
(non à une propriété physique du système : superposition possible). Sa fonction d’onde est
‘indépendante’ [au sens physique] de l’observateur. Le résultat d’une mesure (phénomène
dit de réduction du paquet d’onde), qui peut subvenir, représente lui une propriété physique
du système. Pour le physicien N.Bohr, « le point de vue de la complémentarité se présente
comme une généralisation rationnelle de l’idée même de causalité », reliant une procédure
de préparation d’un système à une mesure d’observable qui le caractérise – sans usage
explicite d’une réduction du vecteur d’état du système (ni donc de la problématique
superposition d’états ne définissant pas de manière unique et claire son état occupé) vu la
non-séparabilité entre ce système et l’appareillage expérimental. On est encore à ce niveau
dans une description plutôt qu’une interprétation. Néanmoins, l’approche probabiliste
(ondulatoire) et la notion de superposition d’états sont inhérentes au formalisme quantique
afin de pouvoir rendre compte de certaines expériences telles les deux fentes de Young
(phénomène d’interférences notamment destructives, variété de positions possibles sur la
plaque photographique). Reste donc à leur donner un sens, une interprétation satisfaisante
voire intuitive...
\subsection{Les interprétations de l’orthodoxie}
N.Bohr a plutôt une vision réaliste de la MQ au sens philosophique (c.f. I.1) : typiquement
l’état avant la mesure (sans propriété physique connue) revêt une réalité [existence]
indépendante de notre perception donc (pour Bohr) de nous-même (et inaccessible aussi).
W.Heisenberg se situe lui dans l’idéalisme philosophique : « Mais les atomes ou les
particules élémentaires ne sont pas réels ; ils forment un monde de potentialités ou de
possibilités, plutôt qu’un monde de choses et de faits ». Pas de réalité indépendante donc
(pour lui) mais une réalité dans l’observation subjective : « Les sciences de la nature ne
donnent pas simplement une description et une explication de la nature ; elles font partie des
interactions entre la nature et nous-mêmes ; elles décrivent la nature telle qu’elle réagit à
notre méthode pour lui poser des questions » [tiré de ‘Physics and philosophy’].
Pour P.Jordan, tel que cité par J.S.Bell : « les observations ne se contentent pas de perturber
ce qui doit être mesuré, elles le produisent. », ce qui rejoint la création de l’objet par le sujet
dans l’idéalisme allemand.
La vision de E.Schrödinger s’accorde de même avec l’idéalisme : « le monde ne m’est
donné qu’une seule fois, et pas une fois pour son existence et une autre pour sa perception.
Le sujet et l’objet sont un. » (pas de réalité existant de manière indépendante du sujet, et,
relation forte objet-sujet).
